\documentclass{exos}
\usepackage{main}

\title{Activité : Compression d'image \& pixel art}
\author{Seconde}
\date{}

\begin{document}
\maketitle

\section{Pixel Art}

\begin{tcolorbox}
Rendez-vous sur \url{https://www.piskelapp.com/kids/} ou munissez-vous d'une feuille à carreaux. Cela devrait vous permettre de faire du pixel art.
\end{tcolorbox}

\begin{alphaquestions}
\item Prévoyez une grille $8 \times 11$;
\item Utilisez les instructions suivantes pour colorier les bons pixels. Le point en bas à gauche est de coordonnées $(0;0)$, et celui en haut à droite est de coordonnées $(7;10)$.
\begin{itemize}

\item Colorier le pixel de coordonnées $(0;3)$.  
\item Colorier le pixel de coordonnées $(0;4)$.  
\item Colorier le pixel de coordonnées $(0;6)$.  
\item Colorier le pixel de coordonnées $(0;7)$.  
\item Colorier le pixel de coordonnées $(1;0)$.  
\item Colorier le pixel de coordonnées $(1;2)$.  
\item Colorier le pixel de coordonnées $(1;8)$.  
\item Colorier le pixel de coordonnées $(1;10)$.  
\item Colorier le pixel de coordonnées $(2;0)$.  
\item Colorier le pixel de coordonnées $(2;2)$.  
\item Colorier le pixel de coordonnées $(2;3)$.  
\item Colorier le pixel de coordonnées $(2;4)$.  
\item Colorier le pixel de coordonnées $(2;5)$.  
\item Colorier le pixel de coordonnées $(2;6)$.  
\item Colorier le pixel de coordonnées $(2;7)$.  
\item Colorier le pixel de coordonnées $(2;8)$.  
\item Colorier le pixel de coordonnées $(2;10)$.  
\item Colorier le pixel de coordonnées $(3;0)$.  
\item Colorier le pixel de coordonnées $(3;1)$.  
\item Colorier le pixel de coordonnées $(3;2)$.  
\item Colorier le pixel de coordonnées $(3;3)$.  
\item Colorier le pixel de coordonnées $(3;4)$.  
\item Colorier le pixel de coordonnées $(3;5)$.  
\item Colorier le pixel de coordonnées $(3;6)$.  
\item Colorier le pixel de coordonnées $(3;7)$.  
\item Colorier le pixel de coordonnées $(3;8)$.  
\item Colorier le pixel de coordonnées $(3;9)$.  
\item Colorier le pixel de coordonnées $(3;10)$.  
\item Colorier le pixel de coordonnées $(4;1)$. 
\item Colorier le pixel de coordonnées $(4;2)$. 
\item Colorier le pixel de coordonnées $(4;4)$. 
\item Colorier le pixel de coordonnées $(4;5)$. 
\item Colorier le pixel de coordonnées $(4;6)$.
\item Colorier le pixel de coordonnées $(5;2)$.
\item Colorier le pixel de coordonnées $(5;3)$.
\item Colorier le pixel de coordonnées $(5;4)$.
\item Colorier le pixel de coordonnées $(5;5)$.
\item Colorier le pixel de coordonnées $(5;6)$.
\item Colorier le pixel de coordonnées $(5;7)$.
\item Colorier le pixel de coordonnées $(5;8)$.
\item Colorier le pixel de coordonnées $(6;3)$.
\item Colorier le pixel de coordonnées $(6;7)$.
\item Colorier le pixel de coordonnées $(7;2)$.
\item Colorier le pixel de coordonnées $(7;8)$.
\end{itemize}
\end{alphaquestions}

\section{Compression d'image}
Comme vous avez pu le constater, les instructions sont très longues. Il faut donc penser à une façon de les condenser pour que cela soit plus rapide à faire.

\begin{tcolorbox}
Avec l'aide d'un ou d'une camarade, trouvez une façon de dicter des instructions afin de reproduire le même dessin que précédemment \textbf{le plus rapidement possible}.
\end{tcolorbox}
\end{document}